\documentclass[11pt,a4paper]{article}
\usepackage[T1]{fontenc}
\usepackage[utf8]{inputenc}
\usepackage{lmodern,microtype}
\usepackage[margin=24mm,top=23mm,bottom=24mm,headheight=15pt]{geometry}
\usepackage{amsmath,amssymb,amsthm,mathtools,booktabs,array,tabularx,enumitem,xcolor,fancyhdr,listings,tikz,float}
\usetikzlibrary{positioning,calc}
\definecolor{ink}{HTML}{17394A}
\definecolor{accent}{HTML}{126D71}
\definecolor{muted}{HTML}{53656D}
\usepackage[colorlinks=true,linkcolor=accent,citecolor=accent,urlcolor=accent,hypertexnames=false,
pdfauthor={Keio and Shaitetsu},
pdftitle={The Sharp Caporossi--Hansen Gap: All-Order Classes and a Million-Vertex Cutoff}]{hyperref}
\newtheorem{theorem}{Theorem}[section]
\newtheorem{lemma}[theorem]{Lemma}
\newtheorem{proposition}[theorem]{Proposition}
\newtheorem{corollary}[theorem]{Corollary}
\theoremstyle{remark}\newtheorem{remark}[theorem]{Remark}
\DeclareMathOperator{\dist}{dist}
\DeclareMathOperator{\diam}{diam}
\DeclareMathOperator{\supp}{supp}
\newcommand{\E}{\mathbb E}
\newcommand{\R}{\mathbb R}
\newcommand{\doi}[1]{\href{https://doi.org/#1}{doi:#1}}
\setlist{topsep=4pt,itemsep=3pt,parsep=0pt,leftmargin=*}
\setlength{\parindent}{1.1em}
\setlength{\parskip}{2pt}
\setlength{\emergencystretch}{3em}
\setlength{\tabcolsep}{5pt}
\renewcommand{\arraystretch}{1.15}
\pagestyle{fancy}\fancyhf{}
\fancyhead[L]{\small\color{muted}The sharp Caporossi--Hansen gap}
\fancyhead[R]{\small\color{muted}All-order classes and an explicit cutoff}
\fancyfoot[C]{\small\thepage}
\renewcommand{\headrulewidth}{0.3pt}
\raggedbottom
\lstset{language=Python,basicstyle=\ttfamily\scriptsize,columns=fullflexible,
keepspaces=true,showstringspaces=false,breaklines=true,breakatwhitespace=false,
frame=single,rulecolor=\color{muted!25},framesep=5pt,
aboveskip=8pt,belowskip=8pt,tabsize=4}
\begin{document}
\thispagestyle{plain}
\noindent{\sffamily\small\bfseries\color{accent}GRAPH THEORY\quad /\quad PROOF CONTINUATION}\par
\vspace{9pt}
\noindent{\LARGE\bfseries\color{ink}The Sharp Caporossi--Hansen Gap}\par
\vspace{6pt}
\noindent{\large All-Order Classes and a Million-Vertex Cutoff}\par
\vspace{9pt}
\noindent Keio and Shaitetsu\hfill 24 November 2026\par
\vspace{10pt}\hrule\vspace{10pt}
\noindent\textbf{Abstract.}
The sharp additive Caporossi--Hansen assertion is
$R(G)-\mu(G)\ge\sqrt{n-1}-2+2/n$, with the star as its unique equality
graph. We continue an earlier corpus-based proof valid for $n\ge10^{10}$.
Using the full weighted-distance theorem of the fixed VMRC source [227],
a shared budget for degree irregularity, and a two-block probability
measure, we lower the sufficient order to $n\ge10^6$.
We also establish the assertion at every order for graphs of minimum
degree at least two, and, independently of the corpus theorem, for all
trees with at most two vertices of degree at least three.
The latter includes every spider, broom, double star, double broom,
and arbitrary subdivisions of the pendant arms of a two-branching-vertex tree.
All finite parameter checks use exact integer or rational arithmetic,
and their complete programs are embedded in this single source.
\emph{This paper does not prove the unrestricted all-order conjecture.}
It identifies the remaining range and records a concrete obstruction to
a tempting cycle-deletion reduction.

\section{The target and the exact scope of the results}

Graphs in this paper are finite, simple, undirected and connected, with
$n\ge2$ vertices. For $d_v=\deg(v)$, put
\[
 R(G)=\sum_{uv\in E(G)}\frac1{\sqrt{d_ud_v}},\qquad
 M(G)=\sum_{uv\in E(G)}\frac1{\max(d_u,d_v)},\qquad
 \mu(G)=\binom n2^{-1}\sum_{\{u,v\}}\dist_G(u,v).
\]
Let
\[
 \gamma_n=\sqrt{n-1}-2+\frac2n,
 \qquad \mathcal S(G)=R(G)-\mu(G)-\gamma_n.
\]
The sharp additive proposal of Caporossi and Hansen \cite{cap}, also
stated explicitly in Dalf\'o \cite{dalfo}, asks whether
\begin{equation}\label{eq:target}\tag{CH}
 \mathcal S(G)\ge0,
 \qquad\mathcal S(G)=0\Longleftrightarrow G\cong K_{1,n-1}.
\end{equation}
The star has $R=\sqrt{n-1}$ and $\mu=2-2/n$, so the constant is forced.
The weaker conclusion $R\ge\mu$ does not imply \eqref{eq:target}.
For comparison, the published tree theorem of Cygan, Pilipczuk and
\v Skrekovski \cite{cygan} gives $R-\mu\ge\max\{0,\sqrt\ell-2\}$
in terms of the number $\ell$ of leaves. It does not supply the
order-dependent sharp gap required in \eqref{eq:target}.

\begin{theorem}[Results established here]\label{thm:summary}
The inequality \eqref{eq:target}, with its stated equality condition,
holds in each of the following classes:
\begin{enumerate}
\item Every connected graph of order $n\ge10^6$.
\item Every connected graph of minimum degree at least two, at every order;
in this class the inequality is strict.
\item Every tree with at most two branching vertices, at every order,
where a branching vertex means a vertex of degree at least three.
\item Every connected graph of diameter at most two, at every order.
\end{enumerate}
Items 1 and 2 use the corpus weighted-distance theorem stated below.
Items 3 and 4 are proved directly.
\end{theorem}

The previous sufficient thresholds $10^{27}$ in VMRC [227] \cite{vmrc227}
and $10^{10}$ in \emph{Graffiti Arithmetic Transfers} \cite{transfers}
are thresholds of those manuscripts, not thresholds attributed to an
external journal paper. The new $10^6$ threshold improves the immediately
preceding one by a factor of $10^4$.

Taking the exact small-order verification for $2\le n\le9$ recorded in
[227] together with the present theorems, the still-unresolved class in
this continuation is contained in
\begin{equation}\label{eq:remaining}
 10\le n\le999\,999,\qquad \delta(G)=1,\qquad \diam(G)\ge3.
\end{equation}
If such a graph is a tree, it must have at least three branching vertices.
Connected graphs with cycles and leaves also remain in scope. No check
of finitely many sample graphs closes this class; no unrestricted closure
is asserted here.

\section{Inputs and elementary estimates}
\label{sec:inputs}

\subsection{The fixed corpus theorem}

The non-elementary input is the weighted-distance theorem in the
Verifiable Mathematical Reasoning Corpus, manuscript [227], at commit
\begin{center}\small\ttfamily
52c1e4ae1f82bd3fdc8b809b9a94e2e4c6ef4cb3.
\end{center}
For every probability vector $p$ on $V(G)$ it states
\begin{equation}\label{eq:weighted-input}\tag{W}
 \sum_{u,v}p_up_v\dist_G(u,v)+2\sum_vp_v^2\le M(G)+1.
\end{equation}
This is a stated dependency on the corpus proof; it is not represented
as a separately located externally refereed theorem.
The preceding eighteen-conjecture synthesis reproduces that proof
\cite{status}. This paper gives the new deductions in full.
Uniform mass in \eqref{eq:weighted-input} gives
\begin{equation}\label{eq:uniform}\tag{CH-U}
 M\ge\frac{(n-1)\mu-(n-2)}n.
\end{equation}
The million-vertex argument will use nonuniform mass as well; its
dependence is therefore on the full statement \eqref{eq:weighted-input}.
The minimum-degree argument uses only \eqref{eq:uniform}.

\subsection{A degree bridge and the maximum-degree broom bound}

\begin{lemma}\label{lem:basic}
Write $d=\Delta(G)$, $r=\sqrt d$ and $k=n-1-d$. Then $M\ge1$ and
\begin{align}
 n-2R&\le(r-1)(R-M),& (r+1)R&\ge n+(r-1)M,\label{eq:bridge}\\
 \mu&\le B(n,d):=
 \frac{n(n-1)(n+1)-(d-1)(d-2)(3n-2d-3)}{3n(n-1)}
 \nonumber\\[-2pt]
 &=2-\frac2n+
 \frac{k(k+1)(3n-2k-7)}{3n(n-1)}.\label{eq:broom}
\end{align}
In particular, $\mu\le(n+1)/3$.
\end{lemma}

\begin{proof}
Every edge at a vertex of maximum degree contributes exactly $1/d$ to
$M$, so these edges contribute one in total. For $1\le a\le b\le d$,
\[
 (a^{-1/2}-b^{-1/2})^2
 =(\sqrt{b/a}-1)\bigl((ab)^{-1/2}-b^{-1}\bigr).
\]
The sum of the squares over edges is $n-2R$, because
$\sum_{uv\in E}(d_u^{-1}+d_v^{-1})=n$.
The factor $\sqrt{b/a}-1$ is at most $r-1$, proving the bridge.

For the mean-distance bound, extend the complete star at a degree-$d$
vertex to a spanning tree. Starting with this $(d+1)$-vertex star, grow
the rest of the tree one leaf at a time. At a step of resulting order
$t$, sort the new leaf's distances as $b_1\le\cdots\le b_{t-1}$.
A path to distance $j$ contains all smaller positive distances, so
$b_i\le i$. A tree with a degree-$d$ vertex has diameter at most
$t-d+1$, giving $b_i\le t-d+1$. Consequently its new transmission
is at most
\[
 \sum_{i=1}^{t-1}\min(i,t-d+1)
 =\binom t2-\binom{d-1}{2}.
\]
The initial star has Wiener sum $d^2$. Summing these increments for
$t=d+2,\ldots,n$, and dividing by $\binom n2$, yields \eqref{eq:broom}.
The broom formed by extending one arm attains this bound.
\end{proof}

\begin{corollary}\label{cor:diameter-two}
Every connected graph of diameter at most two satisfies \eqref{eq:target},
with equality only on a star.
\end{corollary}
\begin{proof}
Put $s=\sqrt{n-1}$. The bridge and $M\ge1$ yield
\[
 R\ge1+\frac{n-2}{r+1}\ge1+\frac{n-2}{s+1}=s.
\]
If $e=|E(G)|$ and the diameter is at most two, then
$\mu=2-e/\binom n2\le2-2/n$. Equality in the latter estimate forces
$e=n-1$; a tree of diameter at most two is a star. The case $K_2$ is
included.
\end{proof}

\subsection{The exact weighted-tree estimate}

For the concentration step we reproduce the weighted-tree lemma from
the immediately preceding manuscript \cite{transfers}. Its precise cubic
deficit is needed; the weaker asymptotic estimate loses too much.
\begin{lemma}\label{lem:weighted-tree}
Let $T$ have $h$ edges and positive integer vertex weights $w_v$ with total
$n$. Let $X,Y$ be independent samples with probability $w_v/n$ at $v$.
Then
\begin{equation}\label{eq:weighted-tree}
 \mathbb E d_T(X,Y)
 \le \frac h2-\frac{h(h^2-1)}{6n^2}
       -\frac{h}{2n^2}\,\boldsymbol{1}_{\{n+h\ \text{even}\}}.
\end{equation}
For $h\ge1$, equality holds precisely when $T$ is a path, every internal
vertex has weight one, and its two endpoint weights differ by at most one.
For $h=0$ the assertion is the equality $0=0$.
\end{lemma}

\begin{proof}
An edge whose deletion leaves weights $a_e,n-a_e$ contributes
$2a_e(n-a_e)/n^2$ to the expected distance. Let
$s_e=\min(a_e,n-a_e)$.

Fix an integer $a$ with $1\le a\le\lfloor n/2\rfloor$. The edges with
$s_e\ge a$, if any, form a connected subtree $S$: every edge on the path
between two such edges separates an outer side of each of those two cuts,
and both outer sides have weight at least $a$. Suppose $S$ has $m\ge1$
edges and $\ell\ge2$ leaves. At each leaf of $S$, the component pointing
away from its unique incident edge in $S$ has weight at least $a$. These
$\ell$ components are disjoint. The other $m+1-\ell$ vertices of $S$
each have weight at least one and lie outside those components. Consequently
\[
 n\ge \ell a+(m+1-\ell)\ge m+2a-1,
 \qquad m\le n-2a+1.
\]
Order the cut sizes as $s_1\ge\cdots\ge s_h$. Taking $a=s_j$ gives
\[
 s_j\le \left\lfloor\frac{n-j+1}{2}\right\rfloor
       \quad(1\le j\le h).
\]
Since $s(n-s)$ increases on $0\le s\le n/2$, the sum of the cut products
is at most the sum of the $h$ largest cut products of the unit-weight path
on $n$ vertices. The corresponding deficits from $n/2$ are the first $h$
entries of
\[
 0,1,1,2,2,\ldots\quad(n\text{ even}),
 \qquad
 \tfrac12,\tfrac12,\tfrac32,\tfrac32,\ldots\quad(n\text{ odd}).
\]
Their squared sum equals
\[
 \frac{h(h^2-1)}{12}
   +\frac h4\,\boldsymbol{1}_{\{n+h\ \text{even}\}}.
\]
Substitute this into
\[
 \mathbb E d_T(X,Y)
   =\frac h2-\frac2{n^2}\sum_e(n/2-s_e)^2
\]
to obtain the bound.

The path in the statement attains exactly these central cut sizes, so the
bound is sharp. For completeness, equality also forces that form. Every
ordered cut size must attain its upper bound, since the integer cut-product
function is strictly increasing up to $\lfloor n/2\rfloor$. If $n$ is even,
there is a central edge cutting the total weight equally. The next required
cut sizes, one on each side, are $n/2-1$, then $n/2-2$, and so on. To obtain
a component of weight one less within a side, the intervening vertex must
have weight one and must have just one continuing branch. This forces two
chains from the central edge, of lengths differing by at most one. If $n$
is odd and $h\ge2$, the two cuts of smaller-side weight $(n-1)/2$ isolate
two disjoint components, leaving a single intervening vertex of weight one;
the same argument then forces the two chains. At the final unpaired cut
either chain may be extended. All edges are exhausted, and the remaining
endpoint weights are the two integers nearest $(n-h+1)/2$. The case $h=1$
states directly that the two endpoint weights are as equal as possible.
\end{proof}


\section{A cutoff of one million from the full weighted-distance theorem}
\label{sec:ch-million}
This section uses the full input \eqref{eq:weighted-input}, including
nonuniform probability mass, in addition to Lemma~\ref{lem:basic}.
For an edge whose endpoint degrees are $a\le b$, write
\begin{equation}\label{eq:chm-edge-identity}
 X_e=\frac1{\sqrt{ab}}-\frac1b,\quad
 S_e=(a^{-1/2}-b^{-1/2})^2
     =(\sqrt{b/a}-1)X_e,\quad
 \sum_eS_e=n-2R.
\end{equation}
\subsection{A two-block use of the weighted theorem}
\begin{lemma}[Two-block estimate]\label{lem:chm-two-block}
Assume \eqref{eq:weighted-input}, and suppose $k=n-1-\Delta\ge1$ and
$\Delta+1\ge k$. Then
\begin{equation}\label{eq:chm-block}
 \mu\le\lambda M+B,\qquad
 \lambda=\frac{4(\Delta+1)(n-\Delta-1)}{n(n-1)},\qquad
 B=\frac{2\Delta}{n-1}.
\end{equation}
\end{lemma}

\begin{proof}
Choose a maximum-degree vertex $v$, and partition $V(G)$ into
$A=N[v]$ and $B_0=V(G)\setminus A$, of orders $a=\Delta+1$ and $k$.
Define the symmetric kernel
$L_{uv}=d_G(u,v)+2\boldsymbol{1}_{\{u=v\}}$, and let $x$ and $z$ be
its averages on $A\times A$ and $B_0\times B_0$, respectively, and
$y$ its average on $A\times B_0$. All pairs in $A$ have distance at
most two, so $x\le2$, whereas $z\ge0$.

Assign total mass $1/2$ uniformly to each block. Inequality
\eqref{eq:weighted-input} gives $(x+2y+z)/4\le M+1$. Put $p=k/n\le1/2$.
For the uniform probability vector on all of $V(G)$, the kernel average
is
\begin{align*}
 \frac{(n-1)\mu+2}{n}
 &=(1-p)^2x+2p(1-p)y+p^2z\\
 &=p(1-p)(x+2y+z)
      +(1-p)(1-2p)x-p(1-2p)z\\
 &\le4p(1-p)(M+1)+2(1-p)(1-2p)\\
 &=\frac{4ak}{n^2}M+\frac{2a}{n}.
\end{align*}
Rearrangement gives \eqref{eq:chm-block}.
\end{proof}

\subsection{A degree split that shares the same gap budget}

\begin{lemma}\label{lem:chm-concentration}
Let $2\le K\le\Delta$ be an integer and put
$t=\sqrt{K-1}$, $r=\sqrt\Delta$, $X=R-M$, and
\[
 q=\frac{t-1}{r-1},\qquad 0\le q<1.
\]
There exist integers $c\ge1$, $b\ge0$, and $0\le h\le n-1$ such that
\begin{align}
 &c\le\frac n{K+1},\qquad b\ge h-c+1,\label{eq:chm-concentration-counts}\\
 &R\ge\frac{(1-q)b+qn-(t-1)X}{2},\label{eq:chm-refined-split}\\
 &\mu\le\frac n{n-1}
       \left(\frac h2-\frac{h^3-h}{6n^2}+7c\right).
       \label{eq:chm-quotient}
\end{align}
The last estimate uses Lemma~\ref{lem:weighted-tree}.
\end{lemma}

\begin{proof}
Call a vertex high if its degree is at least $K$. Choose a maximal set
of $c$ high vertices at pairwise distance at least three. Their closed
radius-one balls are disjoint, which gives $c(K+1)\le n$. Let $A$ be
the union of their closed radius-three balls, and let $b=n-|A|$.
Every high vertex is within distance two of a chosen vertex, so both
ends of any edge incident to a high vertex belong to $A$.

Split the edge quantities in \eqref{eq:chm-edge-identity} into sums
$(S_H,X_H)$ over edges incident to a high vertex and $(S_L,X_L)$ over
the remaining edges. Then
\[
 S_H\le(r-1)X_H,\qquad S_H\le |A|=n-b,\qquad
 S_L\le(t-1)X_L.
\]
For the middle inequality, use
$S_e\le1/d_u+1/d_v$ and sum the incident fractions at each vertex
of $A$. The other two follow from the same identity, since both
endpoint degrees of a low edge are at most $K-1$. In particular,
$X_H\ge S_H/(r-1)$, and hence
\begin{align*}
 n-2R=S_H+S_L
 &\le S_H+(t-1)(X-X_H)\\
 &\le(t-1)X+(1-q)S_H\\
 &\le(t-1)X+(1-q)(n-b).
\end{align*}
This proves \eqref{eq:chm-refined-split}.

Contract each component of the induced graph $G[A]$ to a vertex whose
positive integer weight is its order, and give every vertex outside
$A$ weight one. If there are $c_A$ contracted components, then
$1\le c_A\le c$. Choose a spanning tree of the connected quotient;
it has $h=b+c_A-1$ edges, so $0\le h\le n-1$ and $b\ge h-c+1$.

A component containing $c_i$ chosen balls has internal diameter at
most $7c_i-1$. To see this, connect its balls in the auxiliary graph
whose edges indicate intersection or an edge between the balls. A
chain involving at most $c_i$ balls needs at most $6c_i$ edges within
balls and $c_i-1$ edges between balls. Every simple quotient-tree path
therefore lifts to a walk in $G$ with at most $7c$ additional edges.
The same estimate covers two original endpoints in the same
contracted component.

For two independent uniform original vertices, the projected
vertices have probabilities equal to their weights divided by $n$.
Lemma~\ref{lem:weighted-tree}, with its nonnegative parity correction
discarded, bounds their mean tree distance by
$h/2-(h^3-h)/(6n^2)$. Their original mean distance is
$(n-1)\mu/n$. Adding the lifting cost proves
\eqref{eq:chm-quotient}.
\end{proof}

\subsection{The explicit eventual theorem}

\begin{theorem}\label{thm:ch-million}
Let $G$ be a finite simple connected graph of order $n\ge10^6$
satisfying the full weighted-distance inequality \eqref{eq:weighted-input}.
Then
\begin{equation}\label{eq:chm-conclusion}
 R(G)-\mu(G)\ge\sqrt{n-1}-2+\frac2n,
\end{equation}
with equality precisely for $K_{1,n-1}$. Consequently, the corpus
weighted-distance theorem proves the Caporossi--Hansen additive
bound, including its unique equality graph, for every $n\ge10^6$.
\end{theorem}

\begin{proof}
Write
\[
 N=10^6,\quad s=\sqrt{n-1},\quad r=\sqrt\Delta,\quad
 g=R-\mu,\quad \gamma=s-2+\frac2n,\quad X=R-M.
\]
We show that $g\le\gamma$ forces $G$ to be a star. By the uniform
specialization \eqref{eq:uniform},
\begin{equation}\label{eq:chm-gap-budget}
 0\le X\le g+1+\frac{\mu-2}{n}
       \le s-1+\frac\mu n<s<\sqrt n.
\end{equation}
Here $\mu<n$ suffices for the strict inequality. Choose
\[
 K=\lceil n/175\rceil,\qquad t=\sqrt{K-1}<\sqrt{n/175}.
\]
There are four cases.

\paragraph{Case 1: $\Delta<K$.}
In this case $r\le t<\sqrt n/10$. Equations
\eqref{eq:bridge} and \eqref{eq:chm-gap-budget} yield
$n-2R<n/10$, so $R>9n/20$. Thus
\[
 g>\frac{9n}{20}-\frac{n+1}{3}
   =\frac{7n}{60}-\frac13>\sqrt n>\gamma
 \qquad(n\ge N),
\]
a contradiction.

\paragraph{Case 2: $K\le\Delta\le27n/50$.}
Use $q,c,b,h$ from Lemma~\ref{lem:chm-concentration}. In particular,
$1\le c\le175$. Combining the refined split with
$R=\mu+g$ and the first upper bound on $X$ in
\eqref{eq:chm-gap-budget} gives
\begin{equation}\label{eq:chm-mu-lower}
 \left(2+\frac{t-1}{n}\right)\mu
 \ge qn+(1-q)(h-c+1)
       -(t+1)\gamma-(t-1)\left(1-\frac2n\right).
\end{equation}
Here we used $g\le\gamma$ and $1-q\ge0$.

Put
\[
 x=\frac hn\in[0,1],\quad z=\frac1{\sqrt n}\le\frac1{1000},
 \quad\tau=\frac t{\sqrt n},\quad
 \alpha=\frac{2+\tau z-z^2}{1-z^2}.
\]
Substitute \eqref{eq:chm-quotient} into
\eqref{eq:chm-mu-lower} and divide by $n$. We obtain
\begin{align}
 0\le{}&\alpha\left(\frac x2-\frac{x^3}{6}
                         +\frac{x}{6n^2}+\frac{7c}{n}\right)
             -(1-q)x-q+\frac{(1-q)(c-1)}n\nonumber\\
       &+(\tau+z)\sqrt{1-z^2}-\tau z-3z^2+4z^4.
       \label{eq:chm-cubic-exact}
\end{align}
The last line is exactly
$[(t+1)\gamma+(t-1)(1-2/n)]/n$.

We now bound this expression by a rational cubic. Set
\[
 \tau_- =\frac{3779}{50000},\quad
 \tau_+ =\frac{189}{2500},\quad
 \rho_+ =\frac{147}{200},\quad
 q_0=\frac{127}{1250},\quad
 \alpha_0=\frac{20001}{10000}.
\]
The following rational comparisons hold:
\begin{align}
 &\tau_-^2<\frac1{175}-\frac1N
                \le\tau^2<\frac1{175}<\tau_+^2,
                \qquad\frac{27}{50}<\rho_+^2,
                \label{eq:chm-tau-bounds}\\
 &q\ge\frac{\tau_- -1/1000}{\rho_+ -1/1000}>q_0,
                \qquad
 \alpha\le
 \frac{2+\tau_+/1000-1/N}{1-1/N}<\alpha_0.
                \label{eq:chm-q-alpha-bounds}
\end{align}
For the lower bound on $q$, use $r/\sqrt n\le\sqrt{27/50}<\rho_+$
and observe that $(\tau_- -z)/(\rho_+ -z)$ decreases with $z$.
For the bound on $\alpha$, use
$\alpha=2+(\tau z+z^2)/(1-z^2)$, which increases with both
$\tau$ and $z$ in the ranges under consideration.

For fixed $x\in[0,1]$, $c\ge1$, and $n$, the first line of
\eqref{eq:chm-cubic-exact} decreases with $q$, since its derivative
is $x-1-(c-1)/n\le0$. It increases with $\alpha$ and $c$.
We may therefore replace $q,\alpha,c$ by $q_0,\alpha_0,175$,
respectively. Also
\begin{align*}
 (\tau+z)\sqrt{1-z^2}-\tau z-3z^2+4z^4
 &\le\tau+(1-\tau)z+4z^4\\
 &\le\tau_++\frac{1-\tau_+}{1000}+\frac4{N^2}.
\end{align*}
Using $1/n\le1/N$ and $x\le1$ for the remaining positive error
terms, the right side of \eqref{eq:chm-cubic-exact} is at most
\begin{equation}\label{eq:chm-rational-cubic}
 \beta x-\frac{\alpha_0}{6}x^3+C,
 \qquad \beta=\frac{\alpha_0}{2}-1+q_0=\frac{2033}{20000},
\end{equation}
where
\begin{align}
 C&=-q_0+\tau_++\frac{1-\tau_+}{1000}
       +\frac{174(1-q_0)+1225\alpha_0}{N}
       +\frac{\alpha_0/6+4}{N^2}\nonumber\\
  &=-\frac{449383117913333}{20000000000000000}.
       \label{eq:chm-cubic-constant}
\end{align}
The maximum of $\beta x-\alpha_0x^3/6$ on $x\ge0$ occurs at
$x_* =\sqrt{2\beta/\alpha_0}<319/1000$ and equals
$2\beta x_*/3$. Consequently the expression in
\eqref{eq:chm-rational-cubic} is strictly less than
\begin{equation}\label{eq:chm-negative-certificate}
 \frac23\frac{2033}{20000}\frac{319}{1000}+C
  =-\frac{51095353739999}{60000000000000000}
  <-\frac1{1250}<0.
\end{equation}
This contradicts \eqref{eq:chm-cubic-exact}.

\paragraph{Case 3: $27n/50<\Delta\le13n/20$.}
Here $k=n-1-\Delta\ge1$ and $\Delta+1\ge k$, so
Lemma~\ref{lem:chm-two-block} applies. The function $u(1-u)$
decreases for $u\ge1/2$; with $u=(\Delta+1)/n>27/50$, it gives
\begin{equation}\label{eq:chm-middle-coefficients}
 \lambda<\frac{621n}{625(n-1)}<\frac{497}{500}<1,
 \qquad B\le\frac{13n}{10(n-1)}<\frac32.
\end{equation}
Also $r>\sqrt{27n/50}>700$, and hence
\[
 \frac{r-1}{r+1}>\frac{349}{350}>\frac{497}{500}>\lambda.
\]
The bridge can be written as
$R\ge[n+(r-1)M]/(r+1)$. Combining it with
\eqref{eq:chm-block}, and then using $M\ge1$, yields
\begin{align*}
 g&\ge\frac n{r+1}-B+
           \left(\frac{r-1}{r+1}-\lambda\right)M\\
  &\ge1+\frac{n-2}{r+1}-B-\lambda
   >\frac{n-2}{r+1}-\frac32.
\end{align*}
Since $r\le\sqrt{13n/20}<9\sqrt n/10$ and $n\ge N$, we have
$r+1<\sqrt n$. It follows that
\[
 g>\sqrt n-\frac2{\sqrt n}-\frac32
   >\sqrt n-2+\frac2n>\gamma.
\]
The middle strict inequality follows from
$1/2-2/1000-2/N>0$. This is again a contradiction.

\paragraph{Case 4: $\Delta>13n/20$.}
Put $k=n-1-\Delta$. Suppose first that $k\ge1$ and let
$A_r=2+(r-1)/n$. The bridge and \eqref{eq:uniform} give
\[
 (r+1)g\ge n-2\mu-(r-1)\left(1+\frac{\mu-2}{n}\right).
\]
Use $g\le\gamma$ and the identity
$n-r-3=(r+1)(s-2)+(s-1)(s-r)$ to obtain
\[
 (s-1)(s-r)\le\frac{2(r+1)}n+A_r(\mu-2).
\]
Insert the broom bound \eqref{eq:broom}, use
$s-r=k/(s+r)$, and divide by $k$. This forces
\begin{equation}\label{eq:chm-forced}
 \frac{s-1}{s+r}
 \le\frac{2(r-1)(n-1)}{n^2k}
       +A_r\frac{(k+1)(3n-2k-7)}{3n(n-1)}.
\end{equation}
Put $x=k/n$, so $0<x<7/20$, and let $f(x)=x-2x^2/3$.
The function $f$ is increasing on $[0,7/20]$, since $f'(x)=1-4x/3>0$. The two terms on the right satisfy
\begin{align}
 &\frac{2(r-1)(n-1)}{n^2k}\le\frac2{\sqrt n}
              \le\frac1{500},\qquad
 A_r\le\frac{2001}{1000},\label{eq:chm-high-errors}\\
 &\frac{(k+1)(3n-2k-7)}{3n(n-1)}
              \le f(x)+\frac2n\le f(x)+\frac2N.
              \label{eq:chm-high-fraction}
\end{align}
For the last estimate, put $\varepsilon=1/n$. Its left side is
\[
 \frac{(x+\varepsilon)(1-2x/3-7\varepsilon/3)}{1-\varepsilon}
 \le\frac{f(x)+\varepsilon}{1-\varepsilon}
 \le f(x)+2\varepsilon,
\]
because $f(x)\le f(7/20)=161/600$ and
$161/600+2/N<1$.

If $x\le3/10$, the right side of \eqref{eq:chm-forced} is at most
\[
 \frac1{500}+\frac{2001}{1000}
          \left(\frac6{25}+\frac2N\right)
 <\frac{483}{1000},
\]
whereas its left side is at least
$(s-1)/(2s)>499/1000$. If $3/10<x<7/20$, then
$r<\sqrt{7n/10}<21\sqrt n/25$ and $s-1\ge\sqrt n-2$.
The left side of \eqref{eq:chm-forced} is therefore greater than
\[
 \frac{1-2/1000}{1+21/25}
   =\frac{499}{920}>\frac{27}{50},
\]
while its right side is at most
\[
 \frac1{500}+\frac{2001}{1000}
          \left(\frac{161}{600}+\frac2N\right)
 <\frac{539}{1000}<\frac{27}{50}.
\]
Both alternatives contradict \eqref{eq:chm-forced}. Thus $k=0$.

\paragraph{Equality and completion.}
Only the possibility $\Delta=n-1$ remains. Its dominating vertex
gives $M\ge1$; with $r=s$, the bridge gives
\[
 R\ge\frac{n+(s-1)M}{s+1}
       \ge\frac{n+s-1}{s+1}=s.
\]
All pair distances are at most two, so
\[
 \mu=2-\frac{|E(G)|}{\binom n2}\le2-\frac2n.
\]
If any edge beyond the spanning star is present, this last inequality
is strict. Hence $g\ge\gamma$, and $g\le\gamma$ forces the star.
Conversely the star has $R=\sqrt{n-1}$ and
$\mu=2-2/n$, so it attains equality.
\end{proof}

\paragraph{Scope.}
This argument lowers the sufficient order threshold from $10^{10}$
to $10^6$. It relies on the full weighted-distance theorem through
both its uniform specialization and Lemma~\ref{lem:chm-two-block}.
It does not establish the bound for every order below $10^6$.
If combined with the earlier recorded verification for $2\le n\le9$,
the interval still requiring a universal argument is
$10\le n<10^6$. The accompanying rational-arithmetic certificate
checks the fixed numerical comparisons in this proof; it is not
a verification of the graphs in that remaining interval.

% Input file, not a standalone document.
% Requires amsmath/amssymb and theorem/lemma environments.
% Inputs supplied elsewhere: (CH-U), labelled eq:uniform,
% and the sharp maximum-degree broom bound, labelled eq:broom.
% The full standard-library arithmetic certificate is
% research_external/ch_mindegree2_certificate.py.

\section{The complete case of minimum degree at least two}
\label{sec:ch-mindegree2}

This section closes every order for graphs without leaves. Its inputs are
the uniform weighted-distance inequality \eqref{eq:uniform} and the sharp
broom bound \eqref{eq:broom}; the remaining estimates are proved here.
The finite part checks inequalities on invariant ranges, rather than
enumerating or sampling graphs.

\begin{theorem}\label{thm:ch-mindegree2}
Let $G$ be a connected finite simple graph of order $n$ and minimum degree
at least two. Then
\[
 R(G)-\mu(G)>\sqrt{n-1}-2+\frac2n.
\]
Thus any equality graph or counterexample to the unrestricted
Caporossi--Hansen additive statement must have a leaf.
\end{theorem}

Write $d=\Delta(G)$, $e=|E(G)|$, $Q=\binom n2$ and
\[
 t=\sqrt{d/2},\qquad
 B=B(n,d)=
 \frac{n(n-1)(n+1)-(d-1)(d-2)(3n-2d-3)}
      {3n(n-1)}.
\]
Here $n\ge3$ and $2\le d\le n-1$. The two previously established inputs are
\begin{equation}\label{eq:md2-inputs}
 M\ge\frac{(n-1)\mu-(n-2)}n,\qquad \mu\le B.
\end{equation}

\subsection{Degree ratios, edge counts and distance deficits}

\begin{lemma}\label{lem:md2-bounds}
For every upper bound $A\ge\mu$,
\begin{equation}\label{eq:md2-g}
 R-\mu\ge
 g_{n,t}(A):=
 \frac{n-(2+(t-1)/n)A-(t-1)(1-2/n)}{t+1}.
\end{equation}
Moreover,
\begin{align}
 R&\ge L(n,d):=\frac{n\sqrt{2d}}{d+2},\label{eq:md2-dalfo}\\
 R&\ge P(n,d,e):=
 \frac n2-\frac{t-1}{2(t+1)}(e-n),\label{eq:md2-energy}\\
 \mu&\le A(n,d,e):=B(n,d)-\frac{e-n+1}{Q}.\label{eq:md2-edge-distance}
\end{align}
The permitted integer edge counts lie in the interval
\begin{equation}\label{eq:md2-edge-range}
 n+\left\lceil\frac{d-2}{2}\right\rceil
 \le e\le\left\lfloor\frac{nd}{2}\right\rfloor.
\end{equation}
\end{lemma}

\begin{proof}
For an edge with endpoint degrees $2\le a\le b\le d$, the identity
\[
 (a^{-1/2}-b^{-1/2})^2
 =(\sqrt{b/a}-1)\bigl((ab)^{-1/2}-b^{-1}\bigr)
\]
has $\sqrt{b/a}\le t$. Summing and using
$\sum_{uv\in E}(d_u^{-1}+d_v^{-1})=n$ gives
\[
 n-2R\le(t-1)(R-M),\qquad
 (t+1)R\ge n+(t-1)M.
\]
Combining with \eqref{eq:md2-inputs} yields \eqref{eq:md2-g}, since its
right-hand side is decreasing in $A$. The same argument remains valid
if $t$ is replaced by any larger number.

For \eqref{eq:md2-dalfo}, the degree ratio lies between $2/d$ and $d/2$,
and hence
\[
 \frac{\sqrt{ab}}{a+b}\ge\frac{\sqrt{2d}}{d+2}.
\]
Multiply by $a^{-1}+b^{-1}$ and sum over all edges.

For \eqref{eq:md2-energy}, charge each contribution to
\[
 E=n-2R=\sum_{uv\in E}(d_u^{-1/2}-d_v^{-1/2})^2
\]
to an endpoint of larger degree. A vertex $v$ receives at most $d_v$
charges, each no larger than $(2^{-1/2}-d_v^{-1/2})^2$.
Writing $t_v=\sqrt{d_v/2}$ therefore gives
\[
 \begin{aligned}
 E
 &\le\sum_v(t_v-1)^2
 =\sum_v(d_v-2)\frac{t_v-1}{2(t_v+1)}\\
 &\le\frac{t-1}{2(t+1)}\sum_v(d_v-2)
 =\frac{t-1}{t+1}(e-n).
 \end{aligned}
\]
The coefficient is increasing in $t_v$, and $d_v-2\ge0$.
Since $R=(n-E)/2$, this proves \eqref{eq:md2-energy}.

For \eqref{eq:md2-edge-distance}, choose a vertex of degree $d$ and a
spanning tree containing all of its incident edges. This tree exists
because a star is acyclic and extends to a spanning tree, and it has
maximum degree exactly $d$. Every one of the $e-n+1$ omitted edges
corresponds to a distinct vertex pair whose tree distance is at least two
and whose distance in $G$ is one. Thus the Wiener sum decreases by at
least $e-n+1$ from this tree to $G$. Apply the broom bound to the tree.

Finally, the degree sum is at least $d+2(n-1)$ and at most $nd$,
which proves \eqref{eq:md2-edge-range}.
\end{proof}

\subsection{A uniform analytic argument for \texorpdfstring{$n\ge289$}{n >= 289}}

\begin{lemma}\label{lem:md2-one-variable}
For $0\le x\le1$,
\[
 a(x)-\sqrt{x/2}>\frac1{16},
 \qquad a(x)=\frac13+2x^2-\frac43x^3.
\]
\end{lemma}

\begin{proof}
If $0\le x\le1/2$, then $a(x)\ge1/3+(4/3)x^2$.
Put $y=\sqrt{x/2}$. Differentiation, or the corresponding quartic
Young inequality, gives
\[
 y-\frac{16}{3}y^4\le\frac{3\sqrt[3]{3}}{16}
 \qquad(y\ge0);
\]
the maximum occurs at $y=\sqrt[3]{3}/4$. Since
$9\sqrt[3]{3}<13$, as follows from $2187<2197$, we obtain
\[
 a(x)-\sqrt{x/2}
 \ge\frac13-\frac{3\sqrt[3]{3}}{16}
 >\frac13-\frac{13}{48}=\frac1{16}.
\]
If $1/2\le x\le1$, then
\[
 a(x)-\frac{2x+1}{3}
 =\frac23x(1-x)(2x-1)\ge0.
\]
For $z=\sqrt{2x}\ge1$,
\[
 \frac{2x+1}{3}-\sqrt{x/2}
 =\frac16+\frac{(2z-1)(z-1)}6
 \ge\frac16>\frac1{16}.
\]
\end{proof}

\begin{lemma}\label{lem:md2-large}
For all integers $n\ge289$ and $2\le d\le n-1$,
\[
 g_{n,\sqrt{d/2}}\bigl(B(n,d)\bigr)
 >\sqrt{n-1}-2+\frac2n.
\]
\end{lemma}

\begin{proof}
Set $x=d/n$, $s=\sqrt{n-1}$ and
$f(x)=1-3x^2+2x^3$, so that $a(x)=1-\frac23f(x)$.
We will use
\begin{equation}\label{eq:md2-broom-relaxation}
 B\le\frac n3f(x)+2,\qquad B\le\frac{n+1}{3}.
\end{equation}
For the first inequality, its difference multiplied by $3n(n-1)$
is exactly
\[
 n^2(1-x)(2x^2-4x+5)+n(5x+1)-6>0.
\]
The first summand is nonnegative and the last terms equal
$5d+n-6\ge7$. The second inequality follows directly from the formula
for $B$.

Put $\gamma=s-2+2/n$. Direct subtraction in \eqref{eq:md2-g} gives
\[
 \begin{aligned}
 (t+1)\bigl(g_{n,t}(B)-\gamma\bigr)
 &=n-2B+t+3-(t+1)s-\frac{(t-1)B}{n}-\frac4n\\
 &\ge na(x)-ts-s
   +t\left(\frac23-\frac1{3n}\right)-\frac23-\frac{11}{3n}\\
 &\ge n\bigl(a(x)-\sqrt{x/2}\bigr)
       -\sqrt n-\frac23-\frac{11}{3n}\\
 &>\frac n{16}-\sqrt n-\frac23-\frac{11}{3n}.
 \end{aligned}
\]
We used \eqref{eq:md2-broom-relaxation}, $t=\sqrt{nx/2}$,
$s\le\sqrt n$, and Lemma~\ref{lem:md2-one-variable}; the discarded
$t$ term is nonnegative. If $n\ge289$, then $\sqrt n\ge17$ and
\[
 \frac n{16}-\sqrt n
 =\sqrt n\left(\frac{\sqrt n}{16}-1\right)\ge\frac{17}{16}.
\]
The last displayed lower bound is consequently at least
$17/16-2/3-11/867>0$.
\end{proof}

\subsection{An exact finite arithmetic certificate}

For $3\le n\le288$ and $2\le d\le n-1$ there are exactly $41\,041$
integer pairs. The standard-library certificate
\texttt{ch\_mindegree2\_certificate.py}, reproduced in the appendix,
first tests
\begin{equation}\label{eq:md2-base-test}
 \max\{g_{n,t}(B),\,L(n,d)-B\}
 >\sqrt{n-1}-2+\frac2n.
\end{equation}
Only the following thirteen pairs require the edge-sensitive bounds:
\[
 \begin{gathered}
 (20,6),(21,6),(22,6),(23,6),(24,6),(24,7),(25,6),\\
 (25,7),(26,6),(26,7),(27,7),(28,7),(29,7).
 \end{gathered}
\]
For them, the certificate tests all $699$ integer edge counts in
\eqref{eq:md2-edge-range}, putting $A=A(n,d,e)$ and checking
\begin{equation}\label{eq:md2-edge-test}
 \max\{g_{n,t}(A),\,P(n,d,e)-A,\,L(n,d)-A\}
 >\sqrt{n-1}-2+\frac2n.
\end{equation}
All tests succeed.

For completeness, the decisions use only rational lower bounds.
Let $H=10^6$ and compute by integer square roots
\[
 u=\lfloor H\sqrt{2d}\rfloor,\qquad
 v=\lfloor H\sqrt{n-1}\rfloor.
\]
Use
\[
 t_+=\frac{u+1}{2H},\qquad
 \gamma_+=\frac{v+1}{H}-2+\frac2n,\qquad
 L_-=\frac{nu}{(d+2)H}.
\]
The bridge is valid with $t_+\ge t$. In \eqref{eq:md2-energy},
the coefficient $(t-1)/(2(t+1))$ is increasing, so substitution of
$t_+$ also gives a valid lower bound for $R$. Replacing the target by
$\gamma_+$ and $L$ by $L_-$ therefore produces rational lower bounds
for every tested margin. Python's \texttt{Fraction} and \texttt{isqrt}
perform all decisions exactly; no floating-point computation enters.

The smallest certified margin among the executed tests is
\[
 \frac{23\,130\,813\,383}{15\,982\,498\,350\,000}>0,
\]
at $(n,d)=(27,6)$ in the first stage.
These tests exhaust necessary invariant ranges; they do not presume
that every tested parameter tuple is realizable by a graph.

\begin{proof}[Proof of Theorem~\ref{thm:ch-mindegree2}]
For $n\ge289$, combine \eqref{eq:md2-g} with
Lemma~\ref{lem:md2-large}. For $3\le n\le288$, use the exhaustive
rational certificate \eqref{eq:md2-base-test}--\eqref{eq:md2-edge-test},
whose bounds apply by Lemma~\ref{lem:md2-bounds}. Every resulting
inequality is strict.
\end{proof}

\section{All trees with at most two branching vertices}
\label{sec:ch-two-branch}
A branching vertex of a tree means a vertex of degree at least three.
Write
\[
 \Gamma_n=\sqrt{n-1}-2+\frac2n,\qquad
 \mathcal S(T)=R(T)-\mu(T)-\Gamma_n.
\]
The arguments in this section use only elementary tree identities and a
fixed finite collection of certified integer comparisons. In particular,
they do not use the corpus weighted-distance theorem.

\begin{theorem}\label{thm:ch-two-branch}
Every tree $T$ of order $n\ge2$ with at most two branching vertices satisfies
\[
 R(T)-\mu(T)\ge\sqrt{n-1}-2+\frac2n.
\]
Equality holds if and only if $T$ is a star. Thus the assertion includes all
spiders and all trees obtained by arbitrarily subdividing the pendant arms
and the central path of a double star.
\end{theorem}

\subsection{Brooms and spiders}
A broom is obtained from a star of center degree $d\ge2$ by extending one
arm with $k\ge0$ additional vertices. Its order is $n=d+k+1$. Put
\[
 c=1-\frac1{\sqrt2},\qquad
 A(d)=\frac1{\sqrt2}-\frac12-\frac c{\sqrt d}.
\]
Then $A(d)\ge0$ for $d\ge2$, with equality exactly at $d=2$.

\begin{lemma}\label{lem:ch-all-brooms}
Every nonstar broom has strictly positive CH slack $\mathcal S$.
\end{lemma}
\begin{proof}
For $k\ge1$, summing the edge weights and the tree cut contributions to
the Wiener index gives
\begin{align}
 R&=\frac{k-1}{2}+\frac{d-1}{\sqrt d}
        +\frac1{\sqrt{2d}}+\frac1{\sqrt2},\label{eq:ch-broom-R}\\
 \mu&=2-\frac2n+E,\qquad
 E=\frac{k(k+1)(3n-2k-7)}{3n(n-1)}.
 \label{eq:ch-broom-E}
\end{align}
Thus
\[
 \mathcal S=\frac k2+A(d)
       -\bigl(\sqrt{d+k}-\sqrt d\bigr)-E.
\]
Set $z=k+1$. Completing a square yields
\[
 \frac Ek=\frac{z(3n-5-2z)}{3n(n-1)}
 \le\frac{(3n-5)^2}{24n(n-1)}<\frac38,
\]
because $(3n-5)^2<9n(n-1)$ for $n\ge2$. Consequently
\begin{equation}\label{eq:ch-broom-simple}
 \mathcal S>\frac k8+A(d)
       -\bigl(\sqrt{d+k}-\sqrt d\bigr).
\end{equation}
If $d\ge16$, the radical difference is
$k/(\sqrt{d+k}+\sqrt d)<k/8$. If $k\ge64$, it is smaller than
$\sqrt k\le k/8$. Either condition gives strict positivity.

The remaining rectangle is
\[
 2\le d\le15,\qquad1\le k\le63.
\]
It has 882 points. The exact-arithmetic verifier supplied with this paper
checks \eqref{eq:ch-broom-R}--\eqref{eq:ch-broom-E} at every point.
Its minimum certified lower bound is
\[
 \frac{4304594327}{375000000000}>0,
\]
attained by the certificate at $d=2,k=3$, namely $P_6$. These are finite
parameter comparisons, not sampled graphs. More explicitly, with
$Q=10^{12}$ each term $1/\sqrt j$ is bounded below by
$\lfloor Q/\sqrt j\rfloor/Q$, computed as
$\operatorname{isqrt}(Q^2\mathbin{//}j)/Q$, and
$\sqrt{n-1}$ is bounded above by its integer-square-root ceiling at scale
$Q$. After multiplying by $3n(n-1)Q$, every test is an integer inequality.
The cases $k=0$ are stars and attain equality; $K_2$ is handled directly.
\end{proof}

\begin{lemma}\label{lem:ch-all-spiders}
Every nonstar spider has strictly positive CH slack.
\end{lemma}
\begin{proof}
Let the spider have center degree $d\ge2$, and let $b$ count its arms
of length at least two. Directly summing the three edge types gives
\begin{equation}\label{eq:ch-spider-R}
 R=\frac{n-1-d}{2}+\sqrt d+bA(d).
\end{equation}
For a nonstar, $b\ge1$, so this is at least the Randi\'c index of the
broom having the same $n,d$.

The mean distance is at most that of the same broom. For a direct tree
proof, suppose two arm lengths satisfy $a\ge b\ge2$, and move one unit
of length from the shorter arm to the longer one. With $f(t)=t(n-t)$,
the Wiener-index change is
\[
 f(a+1)-f(b)=(a+1-b)(n-a-b-1)\ge0.
\]
The last factor counts vertices in the other arms. The number of long
arms does not increase, so \eqref{eq:ch-spider-R} also shows that $R$
does not increase. Repeat until only one arm has length greater than one.
The resulting broom therefore has no larger $R$ and no smaller $\mu$.
Apply Lemma~\ref{lem:ch-all-brooms}. For $d=2$ these transformations simply
re-root the same path, which is already included.
\end{proof}

\subsection{Every double broom}
Let $D(a,b,L)$ consist of a path of $L\ge1$ edges between two designated
hubs, with $a,b\ge0$ leaves at its two ends. Its order is
$n=a+b+L+1$. When $L\ge2$, write $x=a+1,y=b+1$. Direct calculation gives
\begin{align}
 R&=\frac{L-2}{2}+\sqrt x+\sqrt y
                  -c\left(\frac1{\sqrt x}+\frac1{\sqrt y}\right),
                  \label{eq:ch-db-R}\\
 W&=(a+b)(n-1)+\sum_{i=1}^L(a+i)(n-a-i),\label{eq:ch-db-W}\\
 W-(n-1)^2&=\frac L6
 \bigl(L^2+3L(a+b)-3L+6ab-3(a+b)+2\bigr).
 \label{eq:ch-db-excess}
\end{align}

\begin{lemma}\label{lem:ch-double-broom-reduction}
For fixed $n,L\ge2$, the smallest slack among the integer double brooms
$D(a,b,L)$ occurs either when $a=0$ or $b=0$, or when $|a-b|\le1$.
\end{lemma}
\begin{proof}
Keep $S=x+y$ fixed and put $z=\sqrt{xy}$. Since $ab=z^2-S+1$, the
part of the slack depending on $z$ is
\[
 q_c(z)-Cz^2,\qquad
 q_c(z)=\sqrt{S+2z}\left(1-\frac cz\right),\qquad
 C=\frac{2L}{n(n-1)}>0.
\]
Its first term has derivative
\[
 q_c'(z)=\frac{1+c/z+cS/z^2}{\sqrt{S+2z}}.
\]
The positive numerator decreases and the denominator increases, so $q_c$
is strictly concave. Subtracting $Cz^2$ preserves strict concavity.
For integer $x,y\ge1$ with fixed sum, the extreme values of $z$ are
$\sqrt{S-1}$ and
$\sqrt{\lfloor S/2\rfloor\lceil S/2\rceil}$. A concave function on
the interval is at least the lesser of its endpoint values. These are
exactly the two cases stated. The first is a broom or a path and was
settled above.
\end{proof}

\begin{lemma}\label{lem:ch-balanced-double-brooms}
Every balanced nonstar double broom with $L\ge2$ has positive slack.
\end{lemma}
\begin{proof}
Put $x=d$, $y=e$, where $e=d$ or $d+1$, so $n=L+d+e-1$.
Summing the consecutive path-cut products in \eqref{eq:ch-db-W} gives
\[
 \mu\le\frac{nL}{2(n-1)}
       -\frac{L(L^2-1)}{6n(n-1)}+\frac{2(n-L-1)}n.
\]
When $e=d+1$, an additional nonpositive parity term has been dropped.
Together with \eqref{eq:ch-db-R}, this implies
\begin{align*}
 \mathcal S\ge{}&\sqrt d+\sqrt e
       -c\left(\frac1{\sqrt d}+\frac1{\sqrt e}\right)
       -\sqrt{n-1}-1
       +\frac{L(L^2+9n-13)}{6n(n-1)}.
\end{align*}
Use $e\ge d$, $n\le L+2d$, $L\ge2$ and $c<3/10$.
In particular,
$c(1/\sqrt d+1/\sqrt e)<3/5$ and $9n-13>0$, so
\[
 \mathcal S>
 2\sqrt d-\sqrt{L+2d}-\frac85+
                 \frac{L^3}{6(L+2d)^2}.
\]
Since $-8/5>-7/4$, this gives the convenient strict bound
\begin{equation}\label{eq:ch-db-coarse}
 \mathcal S>
 2\sqrt d-\sqrt{L+2d}-\frac74+
                 \frac{L^3}{6(L+2d)^2}.
\end{equation}
If $d\ge400$, consider three cases. For $L\le d$, the radical
difference is at least $(2-\sqrt3)\sqrt d>\sqrt d/4\ge5$.
For $d<L\le2d$, that difference is nonnegative and the cubic term is
larger than $d/96\ge400/96>7/4$. For $L>2d$, the difference is larger
than $-\sqrt L$ and the cubic term is larger than $L/24$. Here $L>800$.
The function $L/24-\sqrt L-7/4$ is increasing for $L\ge800$ and at 800
is greater than
\[
 \frac{800}{24}-29-\frac74=\frac{31}{12}>0.
\]
The last case also treats every $L\ge800$ with $d\le399$.

It remains to check the fixed rectangle
\[
 1\le d\le399,\qquad e\in\{d,d+1\},\qquad2\le L\le799.
\]
There are 636,804 parameter points. The point $(d,e,L)=(1,1,2)$ is
$P_3$, which is itself a star and has zero slack. The other 636,803
points are strictly positive by the supplied integer verifier, using
\eqref{eq:ch-db-R} and \eqref{eq:ch-db-excess}. The smallest certified
lower bound is again $4304594327/375000000000>0$, at $P_6$.
This closes every balanced case.
\end{proof}

\begin{lemma}\label{lem:ch-adjacent-double-stars}
Every nonstar double broom with $L=1$ has positive slack.
\end{lemma}
\begin{proof}
Now $S=n=x+y$, and with $z=\sqrt{xy}$,
\[
 R=\sqrt{S+2z}\left(1-\frac1z\right)+\frac1z,
 \qquad \mu=2-\frac2n+\frac{2ab}{n(n-1)}.
\]
The variable part of the slack is $q_1(z)+1/z-Cz^2$.
Differentiation gives
\begin{align*}
 z^3(S+2z)^{3/2}\left(q_1(z)+\frac1z\right)''
 ={}&-z^3-3z^2-6Sz-2S^2+2(S+2z)^{3/2}.
\end{align*}
Since $z\le S/2$, for $S\ge8$ the final term is at most
$2(2S)^{3/2}\le2S^2$. Strict concavity follows. The same endpoint
argument reduces the problem to a star or balanced integer hub degrees.

For a balanced case with $n\ge17$, put $u=\sqrt{n-1}\ge4$.
Both hub degrees are at least $(n-1)/2$, so
\[
 R-\sqrt{n-1}>
 (\sqrt2-1)u-\frac{2\sqrt2}{u}
 \ge\frac72\sqrt2-4>\frac12.
\]
The last inequality follows from $\sqrt2>9/7$. On the other hand,
$2ab/[n(n-1)]<1/2$. Thus the balanced slack is positive.
For $4\le n\le16$, all 91 nonstar choices $2\le x\le n-2$
are covered by the exact verifier. Their minimum certified slack is
$46488264409/3000000000000>0$, at $P_4$. At $n=2,3$ the double stars
are stars. This finishes the adjacent-hub case.
\end{proof}

The preceding three lemmas prove the CH inequality for every double broom,
with equality precisely on stars. The finite certificates used here are
summarized below; they form part of the proof and are independently
reproducible from the supplied integer code.
\begin{center}\small
\begin{tabular}{@{}lrl@{}}\hline
Finite comparison set & Nonstar points & Certified minimum slack\\\hline
Brooms, $2\le d\le15$, $1\le k\le63$ &882&
 $4304594327/375000000000$\\
Balanced double brooms, bounded rectangle &636,803&
 $4304594327/375000000000$\\
Adjacent hubs, $4\le n\le16$ &91&
 $46488264409/3000000000000$\\\hline
\end{tabular}
\end{center}

\subsection{Moving arbitrary pendant arms into the central path}
\begin{proof}[Proof of Theorem~\ref{thm:ch-two-branch}]
Paths and spiders were settled above. Suppose that $T$ has exactly two
branching vertices $u,v$, of degrees $d,e\ge3$, joined by a path of length
$L\ge1$. Every component off this path is itself a pendant path.
Consider an arm of length $a\ge2$ at $u$. Let $B$ be the total length
of the other arms at $u$, and let $C$ be the total length of all arms at
$v$. Then
\[
 n=a+B+1+L+C,\qquad B\ge1,\quad C\ge2.
\]
Shorten the selected arm to length one and insert its $a-1$ removed
vertices into the $u$--$v$ path. The order and the two branching degrees
are preserved.

For the Wiener index write $f(t)=t(n-t)$. The old arm contributes
$\sum_{j=1}^a f(j)$. The old central path has consecutive cut sizes
$a+B+1,\ldots,a+B+L$; the new path adds the sizes $B+2,\ldots,B+a$.
All other edge-cut contributions are unchanged. Hence the exact difference is
\begin{align}
 W(T')-W(T)
 &=\sum_{j=2}^a\bigl(f(B+j)-f(j)\bigr)\nonumber\\
 &=B(a-1)(n-B-a-2)
   =B(a-1)(L+C-1)>0.\label{eq:ch-arm-W}
\end{align}
This is a comparison of the completed transformation; an individual
intermediate shortening need not increase $W$.

For $R$, deleting a degree-two vertex from a pendant arm of length at
least three subtracts $1/2$. Inserting it into a central path of length
at least two adds $1/2$. The last arm shortening from length two to one
has the additional decrease $A(d)\ge0$. If the original central path
has length one, the first subdivision instead contributes the additional
change
\[
 \frac1{\sqrt{2d}}+\frac1{\sqrt{2e}}
       -\frac1{\sqrt{de}}-\frac12
 =-\left(\frac1{\sqrt d}-\frac1{\sqrt2}\right)
    \left(\frac1{\sqrt e}-\frac1{\sqrt2}\right)\le0.
\]
For the completed transformation, the exact change is therefore
\[
 R(T')-R(T)=
 \begin{cases}
  -A(d),&L\ge2,\\[2pt]
  -A(d)-\left(\dfrac1{\sqrt d}-\dfrac1{\sqrt2}\right)
          \left(\dfrac1{\sqrt e}-\dfrac1{\sqrt2}\right),&L=1.
 \end{cases}
\]
Consequently $R(T')\le R(T)$ while \eqref{eq:ch-arm-W} gives
$\mu(T')>\mu(T)$.

Apply this operation to every arm longer than one at both hubs. The
result is a double broom of the same order, with no larger $R$ and no
smaller mean distance. Its slack is positive, since two branching
vertices exclude a star. The original slack is therefore positive too.
The only equality graphs in the classes already handled are stars,
which do attain the required value. This proves the theorem.
\end{proof}


\section{Why the remaining class has not been closed}
\label{sec:remaining}

The preceding arguments prove the classes in Theorem~\ref{thm:summary}.
They do not establish a reduction from an arbitrary graph in
\eqref{eq:remaining} to one of those classes. Two different logical steps
would be required: handling trees with three or more branching vertices,
and handling graphs with both cycles and leaves. Neither is supplied by
the small-order verification or by the new asymptotic cutoff.

A natural route would be to remove a nonbridge edge without increasing
$R-\mu$, and repeat until a tree remains. That route is false in this
unrestricted form.

\begin{proposition}\label{prop:deletion-obstruction}
There is a connected graph in which deleting any nonbridge edge strictly
increases $R-\mu$.
\end{proposition}
\begin{proof}
Take $K_6$ and attach two leaves to each core vertex. The graph has
$18$ vertices, and every core degree is seven. Only its core edges are
nonbridges. By symmetry it suffices to delete one such edge $uv$.
The change in the Randi\'c index is
\[
 \Delta R=
 4\left(\frac1{\sqrt6}-\frac1{\sqrt7}\right)
 +8\left(\frac1{\sqrt{42}}-\frac17\right)-\frac17.
\]
The only increased distances are the nine pairs joining
$\{u,\text{its two leaves}\}$ to
$\{v,\text{its two leaves}\}$. Each increases by one, hence
$\Delta\mu=9/\binom{18}{2}=1/17$.
Integer square-root enclosures at scale $10^{12}$, included in the
constant certificate below, give
\[
 \Delta(R-\mu)>\frac{11}{1000}>0.
\]
Thus every allowed edge deletion worsens this objective.
This graph is an obstruction to the proposed reduction, not a
counterexample to \eqref{eq:target}.
\end{proof}

The all-order tree theorem established here uses a specific transformation
valid when there are at most two branching vertices. Extending it to
arbitrary branch configurations requires an additional proved statement.
The most direct extension of the branch-moving rule also fails.

\begin{proposition}\label{prop:branch-obstruction}
There is a tree with more than two branching vertices for which moving
the middle off-spine component to either end branching vertex of the
spine increases both the Wiener index and $R-\mu$.
\end{proposition}
\begin{proof}
Start with the path $v_0,\ldots,v_{14}$. Attach a two-edge pendant path
at each of $v_3$ and $v_{11}$. Attach a vertex $u$ to $v_7$, and attach
$24$ leaves to $u$. This gives $44$ vertices. Extend each of the two
ends $v_0,v_{14}$ by a path of $228$ new edges, giving $n=500$.
Replace $v_7u$ by $v_3u$ or by $v_{11}u$.
In either case, summing tree-cut contributions $a(n-a)$ gives
\[
 \Delta W=400,\qquad
 \Delta\mu=\frac{800}{500\cdot499}.
\]
Only edges at the two changed attachment vertices contribute to
the change of the Randi\'c index, and their exact sum is
\[
 \Delta R=
 \frac{11}{10}+\frac3{2\sqrt2}
 -\frac5{\sqrt6}-\frac1{5\sqrt3}.
\]
The exact certificate constructs this tree, checks both cut sums, and
encloses the radicals to prove
\[
 \Delta(R-\mu)>\frac7{10000}>0.
\]
Thus neither transfer provides the desired monotonic reduction.
Again, this is not a counterexample to \eqref{eq:target}.
\end{proof}

No general reduction is silently assumed. In particular, a finite interval of
orders is still a mathematical gap until every graph in that interval
is covered by an argument or a completed exhaustive certificate.

\section{Verification and provenance}
\label{sec:verification}

Each displayed theorem is accompanied by its proof. The exceptional
parameter boxes in the minimum-degree and two-branching-vertex arguments
are finite \emph{invariant} boxes whose coverage is established before
computation. The exact certificates then check every tuple in those boxes.
They are therefore proof components, not an extrapolation from randomly
chosen graphs. The analytic parts handle all parameters outside the boxes.

The certificates use Python's standard library only. Integer square roots
produce lower or upper bounds for radicals, and all decisions are integer
or rational inequalities. Decimal values printed for readability play no
role in the decisions. The full programs are reproduced below and are
embedded in this TeX file. No external data file, graph generator or
optimization solver is needed to run them.

The source versions are identified in Section~\ref{sec:inputs} and the
bibliography. The exact conclusions are Theorem~\ref{thm:summary}, with
the unclosed class \eqref{eq:remaining}; no publication-priority or
external-referee-acceptance claim is made.

\clearpage
\appendix
\section{Exact constant and obstruction certificate}
\label{app:constants}
The following certificate checks the rational comparisons in the
million-vertex argument and the cycle-deletion obstruction.

\begin{lstlisting}
"""Exact fixed-constant certificate for CH_million_section.tex.

Run with Python 3. Uses fractions.Fraction only: no floating-point arithmetic,
no graph enumeration, and no claim about the remaining orders 10 <= n < 10**6.
Analytic monotonicity reductions are proved in the accompanying TeX section.
"""

from fractions import Fraction as F


def positive(name, value):
    assert value > 0, (name, value)
    print(f"PASS {name}: margin = {value}")


N = 10**6
z0 = F(1, 1000)
tau_lo, tau_hi = F(3779, 50000), F(189, 2500)
rho_hi = F(147, 200)
q0, alpha0 = F(127, 1250), F(20001, 10000)

# Case 1: sqrt(n)/10 is a valid relaxed degree bound; the quadratic
# 7 y**2/60 - y - 1/3 is positive and increasing for y >= 1000.
positive("1/sqrt(175) < 1/10 (squared)", F(1, 100) - F(1, 175))
positive("case 1 endpoint", F(7, 60) * N - 1000 - F(1, 3))
positive("case 1 derivative", F(7, 30) * 1000 - 1)

# Case 2: square-root enclosures and the rational cubic upper bound.
positive("tau lower bound", F(1, 175) - F(1, N) - tau_lo**2)
positive("tau upper bound", tau_hi**2 - F(1, 175))
positive("rho upper bound", rho_hi**2 - F(27, 50))
positive("q lower bound", (tau_lo-z0)/(rho_hi-z0) - q0)
positive("alpha upper bound", alpha0 - (2+tau_hi*z0-z0**2)/(1-z0**2))

beta = alpha0/2 - 1 + q0
assert beta == F(2033, 20000)
C = (-q0 + tau_hi + (1-tau_hi)*z0
     + (174*(1-q0) + 1225*alpha0)/N
     + (alpha0/6+4)/N**2)
assert C == -F(449383117913333, 20000000000000000)
positive("cubic maximizer enclosure", F(319, 1000)**2 - 2*beta/alpha0)
cubic_bound = F(2, 3)*beta*F(319, 1000) + C
assert cubic_bound == -F(51095353739999, 60000000000000000)
positive("negative cubic margin beyond 1/1250", -cubic_bound-F(1, 1250))

# Case 3: the relevant rational functions decrease with n >= N.
positive("middle lambda bound", F(497, 500) - F(621*N, 625*(N-1)))
positive("middle B bound", F(3, 2) - F(13*N, 10*(N-1)))
positive("r > 700 (squared)", F(27*N, 50) - 700**2)
positive("positive M coefficient", F(349, 350)-F(497, 500))
positive("r < 9 sqrt(n)/10 (squared)", F(81, 100)-F(13, 20))
positive("r+1 < sqrt(n) endpoint", 1-F(9, 10)-z0)
positive("middle final gap", F(1, 2)-2*z0-F(2, N))

# Case 4: the error bounds are largest at n=N.
positive("fraction relaxation", 1-F(161, 600)-F(2, N))
positive("s > 500 (squared)", F(N-1)-500**2)
upper_low = F(1, 500)+F(2001, 1000)*(F(6, 25)+F(2, N))
positive("high-degree x<=3/10 bound", F(483, 1000)-upper_low)
positive("high-degree first separation", F(499, 1000)-F(483, 1000))
positive("sqrt(7/10) < 21/25 (squared)", F(21, 25)**2-F(7, 10))
lower_high = (1-2*z0)/(1+F(21, 25))
assert lower_high == F(499, 920)
positive("high-degree second lower bound", lower_high-F(27, 50))
upper_high = F(1, 500)+F(2001, 1000)*(F(161, 600)+F(2, N))
positive("high-degree second upper bound", F(539, 1000)-upper_high)
positive("high-degree second separation", F(27, 50)-F(539, 1000))

print("All exact fixed-constant checks passed.")


# The cycle-deletion obstruction of Section 6.
from math import isqrt
from fractions import Fraction as F
scale=10**12
def inv_interval(d):
    a=isqrt(scale*scale//d)
    return F(a,scale),F(a+1,scale)
a6,b6=inv_interval(6)
a7,b7=inv_interval(7)
a42,b42=inv_interval(42)
change_lower=4*a6-4*b7+8*a42-F(9,7)-F(1,17)
assert change_lower>F(11,1000)
print('Cycle deletion increases the gap by more than 11/1000.')

# The 500-vertex branch-transfer obstruction.
def tree_wiener(n,edges):
    adj=[[] for _ in range(n)]
    for u,v in edges:
        adj[u].append(v)
        adj[v].append(u)
    parent=[-1]*n
    parent[0]=0
    order=[0]
    for u in order:
        for v in adj[u]:
            if parent[v]<0:
                parent[v]=u
                order.append(v)
    assert len(edges)==n-1 and len(order)==n
    size=[1]*n
    for v in reversed(order[1:]):
        size[parent[v]]+=size[v]
    return sum(size[v]*(n-size[v]) for v in order[1:])
edges=[(i,i+1) for i in range(14)]
edges += [(3,15),(15,16),(11,17),(17,18),(7,19)]
edges += [(19,j) for j in range(20,44)]
next_vertex=44
for end in (0,14):
    previous=end
    for step in range(228):
        edges.append((previous,next_vertex))
        previous=next_vertex
        next_vertex+=1
assert next_vertex==500
old_w=tree_wiener(500,edges)
from collections import Counter
def edge_products(n,edge_list):
    deg=[0]*n
    for u,v in edge_list:
        deg[u]+=1
        deg[v]+=1
    return Counter(deg[u]*deg[v] for u,v in edge_list)
old_products=edge_products(500,edges)
for end in (3,11):
    new_edges=[e for e in edges if e!=(7,19)]+[(end,19)]
    assert tree_wiener(500,new_edges)-old_w==400
    new_products=edge_products(500,new_edges)
    diff={p:new_products[p]-old_products[p]
          for p in set(old_products)|set(new_products)
          if new_products[p]!=old_products[p]}
    assert diff=={4:2,8:3,6:-5,75:-1,100:1}
a2,b2=inv_interval(2)
a3,b3=inv_interval(3)
delta_R_lower=F(11,10)+F(3,2)*a2-5*b6-F(1,5)*b3
delta_gap_lower=delta_R_lower-F(800,500*499)
assert delta_gap_lower>F(7,10000)
print('Both branch transfers increase the gap by more than 7/10000.')

\end{lstlisting}
\section{Minimum-degree certificate}\label{app:md2}
\begin{lstlisting}
#!/usr/bin/env python3
"""Exact finite certificate for the CH additive inequality when delta >= 2.

The analytic argument in ch_mindegree2_proof.md handles every n >= 289.
This program checks n = 3,...,288 and all integer Delta = 2,...,n-1.
It uses rational arithmetic only; square roots are enclosed by integer sqrt.
It checks bounds on graph invariants, not a sample or enumeration of graphs.
Inputs: CH-U, the maximum-degree broom bound, and the elementary bounds
proved in the accompanying argument. No numerical solver is required.
"""
from fractions import Fraction as Q
from math import isqrt
import json


SCALE = 10**6


def broom_mean(n, d):
    return Q(n*(n-1)*(n+1) - (d-1)*(d-2)*(3*n-2*d-3),
             3*n*(n-1))


def bounds(n, d):
    # sqrt(2*d) >= a/SCALE and sqrt(d/2) < (a+1)/(2*SCALE).
    a = isqrt(2*d*SCALE*SCALE)
    b = isqrt((n-1)*SCALE*SCALE)
    r_upper = Q(a+1, 2*SCALE)
    s_upper = Q(b+1, SCALE)
    target_upper = s_upper - 2 + Q(2, n)
    dalfo_lower = Q(n*a, (d+2)*SCALE)
    return r_upper, target_upper, dalfo_lower


def bridge_gap(n, A, r_upper, target_upper):
    # A is an upper bound on mu. The ratio parameter r_upper exceeds
    # sqrt(Delta/2), so the edge identity directly justifies this bridge.
    numerator = n - (2+(r_upper-1)/n)*A - (r_upper-1)*(1-Q(2,n))
    return numerator/(r_upper+1) - target_upper


def main():
    base_cases = 0
    edge_cases = 0
    exceptions = []
    best = None
    for n in range(3, 289):
        for d in range(2, n):
            base_cases += 1
            B = broom_mean(n, d)
            r, target, dalfo = bounds(n, d)
            first = max(bridge_gap(n, B, r, target), dalfo-B-target)
            if first > 0:
                if best is None or first < best[0]:
                    best = (first, (n,d,None))
                continue

            exceptions.append((n,d))
            q = Q(n*(n-1), 2)
            # Degree sum: 2e >= 2(n-1)+d. Simple graph: 2e <= n*d.
            lower_e = n + (d-1)//2
            upper_e = n*d//2
            for e in range(lower_e, upper_e+1):
                edge_cases += 1
                # An n-1 edge spanning tree keeps a full maximum-degree
                # star. Each of the e-n+1 omitted edges improves its own
                # endpoint-pair distance by at least one.
                A = B-Q(e-n+1, q)
                energy_R = Q(n,2)-(r-1)/(2*(r+1))*(e-n)
                gap = max(bridge_gap(n,A,r,target),
                          energy_R-A-target,
                          dalfo-A-target)
                assert gap > 0, (n,d,e,gap)
                if best is None or gap < best[0]:
                    best = (gap, (n,d,e))

    expected = [(20,6),(21,6),(22,6),(23,6),(24,6),(24,7),
                (25,6),(25,7),(26,6),(26,7),(27,7),(28,7),(29,7)]
    assert exceptions == expected
    assert base_cases == 41041
    # Constants used in the analytic n >= 289 proof.
    assert 9**3*3 < 13**3
    assert Q(17,16)-Q(2,3)-Q(11,3*289) > 0
    report = dict(base_cases=base_cases, edge_cases=edge_cases,
                  exceptional_pairs=exceptions,
                  minimum_certified_margin=str(best[0]),
                  minimum_margin_location=best[1],
                  square_root_scale=SCALE,
                  result='Every finite case is strictly positive.')
    print(json.dumps(report, indent=2))


if __name__ == '__main__':
    main()

\end{lstlisting}
\section{Broom and double-broom certificate}\label{app:brooms}
\begin{lstlisting}
"""Finite exact comparisons used with the analytic broom/double-broom proof.
Only integer arithmetic and integer square roots; Python3 standard library.
"""
from math import isqrt
from fractions import Fraction
import json
S=10**12
inv=lambda d:isqrt(S*S//d)
rootup=lambda d:(lambda x:x+(x*x<d*S*S))(isqrt(d*S*S))

def check_brooms():
    count=0;best=None;arg=None
    for d in range(2,16):
        for k in range(1,64):
            n=d+k+1
            rlo=(d-1)*inv(d)+inv(2*d)+inv(2)+S*(k-1)//2
            den=3*n*(n-1);num=k*(k+1)*(3*n-2*k-7)
            margin=Fraction((rlo-rootup(n-1))*den-S*num,S*den)
            assert margin>0,(d,k,margin)
            if best is None or margin<best:best,arg=margin,(d,k,n)
            count+=1
    return {'count':count,'minimum_certified':str(best),'minimum_decimal':float(best),'arg':arg}

def check_balanced_double_brooms():
    # Analytic inequalities handle min(hubdegree)>=400 or hub-pathlength>=800.
    iv={d:inv(d) for d in range(1,802)}
    rt={d:rootup(d) for d in range(1,1600)}
    count=0;best=None;arg=None
    for d in range(1,400):
        for e in (d,d+1):
            a,b=d-1,e-1
            base=(d-1)*iv[d]+(e-1)*iv[e]+iv[2*d]+iv[2*e]
            for L in range(2,800):
                if (d,e,L)==(1,1,2):continue # P3 is itself a star.
                n=L+d+e-1
                rlo=S*(L-2)//2+base
                num=L*(L*L+3*L*(a+b)-3*L+6*a*b-3*(a+b)+2)
                den=3*n*(n-1)
                top=(rlo-rt[n-1])*den-S*num
                assert top>0,(d,e,L,n,top)
                count+=1
                if best is None or top*best.denominator<best.numerator*S*den:
                    best=Fraction(top,S*den);arg=(d,e,L,n)
    return {'count':count,'minimum_certified':str(best),'minimum_decimal':float(best),'arg':arg}

def check_adjacent_hubs():
    count=0;best=None;arg=None
    # For n>=8, concavity reduces to stars or balanced hub degrees;
    # n>=17 balanced cases are treated analytically. Exhaust the smaller set.
    for n in range(4,17):
        for d in range(2,n-1):
            e=n-d;a,b=d-1,e-1
            rlo=a*inv(d)+b*inv(e)+inv(d*e)
            den=n*(n-1);num=2*a*b
            margin=Fraction((rlo-rootup(n-1))*den-S*num,S*den)
            assert margin>0,(n,d,e,margin)
            if best is None or margin<best:best,arg=margin,(n,d,e)
            count+=1
    return {'count':count,'minimum_certified':str(best),'minimum_decimal':float(best),'arg':arg}

if __name__=='__main__':
    result={'brooms':check_brooms(),
            'balanced_double_brooms':check_balanced_double_brooms(),
            'adjacent_hubs':check_adjacent_hubs()}
    with open('CH_agent_broom_certificate_results.json','w') as f:json.dump(result,f,indent=2)
    print(json.dumps(result,indent=2))

\end{lstlisting}

\begin{thebibliography}{9}
\bibitem{cap}
G.~Caporossi and P.~Hansen,
Variable neighborhood search for extremal graphs: 1. The AutoGraphiX system,
\emph{Discrete Mathematics} \textbf{212} (2000), 29--44.
\doi{10.1016/S0012-365X(99)00206-X}.
\bibitem{dalfo}
C.~Dalf\'o, On the Randi\'c index of graphs,
\emph{Discrete Mathematics} \textbf{342} (2019), 2792--2796.
\doi{10.1016/j.disc.2018.08.020}.
\bibitem{cygan}
M.~Cygan, M.~Pilipczuk and R.~\v Skrekovski,
Relation between Randi\'c index and average distance of trees,
\emph{MATCH Communications in Mathematical and in Computer Chemistry}
\textbf{66} (2011), 605--612.
\href{https://users.fmf.uni-lj.si/skreko/Old/MyPapers/01130.pdf}{Author preprint}.
\bibitem{vmrc227}
Verifiable Mathematical Reasoning Corpus, manuscript [227],
\emph{Fajtlowicz Proof Audit}, source at commit
\texttt{52c1e4ae1f82bd3fdc8b809b9a94e2e4c6ef4cb3}.
\href{https://github.com/QEEDDD/Verifiable-Mathematical-Reasoning-Corpus/blob/52c1e4ae1f82bd3fdc8b809b9a94e2e4c6ef4cb3/%5B227%5D%20Fajtlowicz_Proof_Audit.tex}{Frozen source}.
The weighted statement is (W); the earlier sufficient cutoff is in the
section on the Caporossi--Hansen gap.
\bibitem{transfers}
Keio and Shaitetsu, \emph{Graffiti Transfers: A Sharper Extremal Bound,
Arithmetic Counterexamples, and a Log-Rank Reformulation},
research manuscript, 24 November 2026.
\href{https://chatgpt.com/library?file_id=libfile_fde3_f-1194bbf4fa3c819195f961b370223ba2_67c74a3277a88191acd4edb2733e13fb&account_id=personal}{Identified TeX predecessor},
\texttt{Graffiti\_Arithmetic\_Transfers.tex}.
\bibitem{status}
Keio and Shaitetsu, \emph{The Eighteen Conjectures of Fajtlowicz (1988):
Truth Values, Proofs, and Counterexamples}, research synthesis,
24 November 2026.
\href{https://chatgpt.com/library?file_id=libfile_fde3_f-d8451b980f7c8191958608b0992ffc0e_67c74a3277a88191acd4edb2733e13fb&account_id=personal}{Identified PDF synthesis}.
\end{thebibliography}
\end{document}
